\newpage
\begin{center}
  \textbf{\large АННОТАЦИЯ}
\end{center}

В магистерской квалификационной работе рассматривается задача реконструкции
простых бинарных визуальных стимулов размером $6 \times 6$ пикселей по
многоканальным сигналам электроэнцефалографии, зарегистрированным в фазе
воспоминания ранее предъявленного изображения. Задача формулируется как
многовыходное бинарное предсказание 36 пикселей, что позволяет исследовать
локальную структуру визуального образа без использования внешних семантических
латентных пространств.

В работе описан воспроизводимый экспериментальный пайплайн, построенный поверх
ранее подготовленного EEG/EOG-корпуса. Для основной задачи выделены 180
наблюдений случайных стимулов фазы воспоминания. Пайплайн включает подготовку
входных EEG-представлений, извлечение временных, спектральных,
пространственных и локально-паттерновых признаков, построение спектральных
тензоров FFT/Welch, Morlet, Superlet и STFT, обучение классических моделей
машинного обучения, референсной логистической регрессии и компактных
сверточных нейросетевых архитектур. Оценка проводится в двух протоколах:
межсубъектном и двунаправленном межпробном, допускающем совпадение испытуемых,
но разделяющем экспериментальные пробы. Для интерпретации результатов
используются попиксельная balanced accuracy, субъектно-кластерный bootstrap и
парное сравнение моделей с поправкой Holm.

Полученные результаты показывают, что завершенный пайплайн обеспечивает
контролируемую и воспроизводимую проверку задачи реконструкции, однако
качество моделей остается близким к случайному уровню. Лучшие описательные
значения balanced accuracy составили 0.518382 в межсубъектном протоколе и
0.512011 в двунаправленном межпробном протоколе; статистически надежного
превосходства над референсной логистической регрессией не выявлено. Для
контроля работоспособности программного контура дополнительно выполнены
эксперименты на открытых корпусах BNCI2014-001 и BNCI2014-009. Они показывают,
что реализованные процедуры загрузки данных, субъектно-раздельной оценки,
обучения и агрегации метрик способны выявлять сигнал на независимых EEG-задачах,
но не являются внешней валидацией реконструкции случайных изображений. Работа
тем самым фиксирует инженерно воспроизводимую основу для дальнейших
исследований и методологические ограничения реконструкции случайных визуальных
образов из EEG при малом объеме данных.

\onehalfspacing
